\documentclass[11pt,a4paper]{article}

\usepackage[utf8]{inputenc}
\usepackage[french]{babel}
\usepackage[margin=2.5cm]{geometry}
\usepackage{parskip}

\begin{document}

\thispagestyle{empty}

\begin{center}
  {\large\textsc{Université de Mons} --- Faculté des Sciences}\\[0.3em]
  {\normalsize Master en Sciences Informatiques --- Juin 2026}\\[1.2em]
  {\Large\bfseries Reconstruction 3D des mouvements de force athlétique\\
  à partir de vidéos monoculaires}\\[0.8em]
  {\normalsize \textbf{Nicolas Delplanque} \quad --- \quad Directeur : Stéphane Dupont}
\end{center}

\vspace{1em}

\noindent\textbf{Résumé.}
L'analyse biomécanique de la force athlétique (flexion de jambes, développé couché, soulevé de
terre) repose traditionnellement sur l'observation visuelle d'un expert ou sur des
capteurs intrusifs, deux approches limitées par la perte d'information spatiale, les
occlusions liées à l'équipement et le besoin de matériel spécialisé. Ce mémoire propose
un pipeline non invasif capable de reconstruire la pose tridimensionnelle d'un athlète à
partir d'une simple vidéo monoculaire non calibrée, sans marqueur ni système
multicaméras.

L'architecture proposée repose sur quatre modules de l'état de l'art : segmentation
d'instances (SAM3), estimation monoculaire de profondeur (MoGe-2), complétion amodale
des occlusions par modèle de diffusion (TACO) et reconstruction de la pose en 3D (SAM 3D Body).
La contribution centrale réside dans la spécialisation du module de complétion amodale à
la force athlétique et dans l'unification de ces composants au sein d'un pipeline complet.
À cette fin, un jeu de données dédié a été constitué à partir de six sujets, avec génération
synthétique d'occlusions représentatives (disques, racks) et augmentation des données,
puis exploité pour un fine-tuning de TACO par adaptation de faible rang (LoRA).

L'évaluation, conduite sur un sujet absent de l'entraînement, montre une amélioration
moyenne du LPIPS (\textit{Learned Perceptual Image Patch Similarity}) de 35,0\,\% par rapport 
au modèle de base (significativité confirmée par un test de Wilcoxon), le gain étant le plus marqué 
au développé couché, mouvement le plus fortement occulté. Mesuré par le MPJPE 
(\textit{Mean Per-Joint Position Error}), le fine-tuning améliore nettement la reconstruction 3D par 
rapport au modèle de base, avec un gain pouvant atteindre $40,6$\,\% au soulevé de terre. Par rapport à l'image occluse brute, 
la complétion amodale n'apporte toutefois un gain intéressant qu'au développé
couché et peut introduire des artefacts pénalisant l'estimation de la pose au squat et au soulevé de terre.
Une interface web de visualisation complète le travail en calculant automatiquement les indicateurs
biomécaniques pertinents (profondeur de squat, verrouillage articulaire, inclinaison du
buste, asymétrie des hanches). Les principales limites tiennent à la taille restreinte du
jeu de données de tests et à l'usage d'occlusions synthétiques, qui appellent une validation sur un
corpus réel élargi.

Les principales limites tiennent à la taille restreinte du jeu de test (un seul sujet en évaluation), 
à l'usage d'occlusions synthétiques plutôt que d'occlusions réelles, et aux artefacts que la complétion 
amodale peut introduire sur certains mouvements. Plusieurs pistes d'amélioration en découlent : valider le pipeline sur 
un corpus réel et multi-sujets, affiner la complétion pour limiter ces artefacts résiduels, et exploiter 
la cohérence temporelle entre images successives, actuellement absente, afin de stabiliser la reconstruction 3D.

\vspace{1em}

\noindent\textbf{Mots-clés :} reconstruction 3D ; force athlétique ; vidéo monoculaire ;
complétion amodale ; fine-tuning LoRA ; estimation de pose humaine.

\end{document}